\scanpage{223}
% opening, (3)
Cette fonction est déterminée, indépendamment du choix de $C_s$, modulo une expression de la forme
\[
 \sum D_{i\,s}\int_{Z_s}\omega_{i\,s}
 \deleted{=\int_{Z_s}\sum_i D_{i\,s}\omega_{i\,s}}
 \tag*{(3)}
\]

% cancelled NB and the definition
\deleted{(\uline{NB} Le fait que $Z_s$ est un 1-\uline{cycle} permet \unsure{de}}

Supposons que les $D_i$, $\omega_i$ soient tels que l'expression précédente soit nulle, quelle que soit \deleted{\unsure{est}} la famille continue $s\unsure{\longmapsto|\longrightarrow}Z_s$ de 1-cycles sur un ouvert de $S$. \deleted{\unsure{$Z$|$\Sigma$}} \deleted{Alors l'expression} \inserted{On dira} que \deleted{les} $\mu=(D_i,\omega_i)_{i\in I}$ \uline{définit une équation de Picard-Fuchs sur} $X/S$. Alors pour deux sections hol. quelc. $a,b$ de $X/S$, l'expression
\[
 \mu(a,b)=\sum_i D_{i\,s}\int_{a(s)}^{b(s)}\omega_{i\,s}
 \tag*{(4}
\]
est bien déterminée, et est une fonction holomorphe sur $S$. Bien sûr on aura
\[
 \mu(a,b)+\mu(b,c)=\mu(a,c).\tag*{(5)}
\]

% last lines, continuing on the next sheet
Cherchons une condition suffisante maniable, \inserted{sous forme différentielle} \deleted{\unsure{pour}} \deleted{un \unsure{critère} \uline{infinitésimal}, et non plus global,} pour que $\mu$ soit de Pic. Fuchs. Notons
% archive sheet number: 222
